% LƯU Ý: Thêm \usepackage{needspace} vào phần khai báo gói
% trong file do-an.tex, trước \begin{document}.
	% File này dùng gói enumitem (đã có sẵn trong settings hoặc do-an.tex).
	
	% ---------------------------------------------------------------------
	% CHƯƠNG 2 — CƠ SỞ LÝ THUYẾT VÀ CÔNG NGHỆ
	% ---------------------------------------------------------------------
	
	\chapter{Cơ sở lý thuyết và công nghệ}
	\label{ch:cosolythuyet}
	
	% Cấu hình danh sách gạch đầu dòng: thụt vào, dòng đầu thẳng hàng với các dòng sau
	\setlist[itemize,1]{
		label=--,
		leftmargin=1.5cm,
		labelsep=0.5em,
		itemindent=0pt,
		listparindent=0pt,
		parsep=0pt,
		itemsep=4pt,
		topsep=4pt
	}
	
	% ---------------------------------------------------------------------
	% 2.1. KIẾN THỨC NỀN TẢNG
	% ---------------------------------------------------------------------
	
	\Needspace{4\baselineskip}
	\section{Kiến thức nền tảng}
	
	\Needspace{3\baselineskip}
	\subsection{Dị ứng thực phẩm và thành phần gây dị ứng}
	
	Dị ứng thực phẩm là phản ứng của hệ miễn dịch đối với một hoặc nhiều
	thành phần có trong thực phẩm. Tùy thuộc vào cơ địa và mức độ nhạy cảm
	của từng người, phản ứng dị ứng có thể biểu hiện ở nhiều mức độ khác
	nhau. Một số trường hợp có thể xuất hiện các triệu chứng như ngứa,
	phát ban, sưng môi, khó thở hoặc những phản ứng nghiêm trọng hơn.
	
	Trong quá trình lựa chọn thực phẩm, việc kiểm tra thành phần và thông
	tin cảnh báo trên bao bì có ý nghĩa quan trọng đối với người có tiền
	sử dị ứng. Tuy nhiên, thông tin về dị nguyên không phải lúc nào cũng
	được trình bày dưới cùng một cách diễn đạt. Một chất có thể xuất hiện
	dưới tên nguyên liệu, tên khoa học hoặc tên gọi thông dụng khác nhau.
	
	Đối với hệ thống Allergy Warning, dị nguyên được xem xét thông qua
	hai nguồn thông tin chính:
	
	\begin{itemize}
		\item Danh sách thành phần: Các nguyên liệu được ghi trên nhãn sản phẩm. Hệ thống phân tích nội dung này để nhận diện những thành phần có khả năng liên quan đến dị ứng.
		\item Thông tin cảnh báo riêng trên nhãn: Các nội dung do nhà sản xuất công bố, chẳng hạn như ``Contains'' (Chứa) hoặc ``May contain'' (Có thể chứa). Đây là thông tin cảnh báo được trình bày riêng trên bao bì sản phẩm.
	\end{itemize}
	
	Hai nguồn thông tin trên có ý nghĩa khác nhau. Một chất có thể xuất
	hiện trong danh sách thành phần nhưng không được trình bày lại trong
	mục cảnh báo riêng của nhà sản xuất. Vì vậy, không thể chỉ dựa vào
	mục cảnh báo riêng để kết luận sản phẩm không chứa chất gây dị ứng.
	
	Trong phạm vi đồ án, hệ thống sử dụng thông tin thành phần và dữ liệu
	dị nguyên để phân tích sản phẩm. Kết quả được đối chiếu với hồ sơ dị
	ứng mà người dùng đã khai báo nhằm xác định những thành phần cần
	được chú ý.
	
	\Needspace{3\baselineskip}
	\subsection{Nhãn thành phần thực phẩm và thông tin cảnh báo dị nguyên}
	
	Nhãn thực phẩm là nguồn cung cấp thông tin về tên sản phẩm, thương
	hiệu, thành phần nguyên liệu và các nội dung cảnh báo liên quan.
	Thông tin này thường được trình bày dưới dạng văn bản trên bao bì.
	
	Trong quá trình phân tích nhãn, cần phân biệt một số loại thông tin:
	\begin{itemize}
		\item Thành phần nguyên liệu: Liệt kê các chất được sử dụng để sản xuất thực phẩm, ví dụ như bột mì, sữa, dầu thực vật hoặc các loại hạt.
		\item Dị nguyên khai báo riêng: Những chất được nhà sản xuất ghi rõ trong phần cảnh báo dị nguyên trên bao bì.
		\item Cảnh báo có thể chứa vết: Thông tin cảnh báo về khả năng sản phẩm có thể chứa một lượng nhỏ chất gây dị ứng do quá trình sản xuất hoặc đóng gói.
	\end{itemize}

	Việc phân biệt các loại thông tin này giúp hạn chế nhầm lẫn giữa
	thành phần thực tế được liệt kê và nội dung cảnh báo do nhà sản xuất
	công bố.
	
	Đối với Allergy Warning, nội dung trên nhãn được sử dụng làm dữ liệu
	đầu vào cho quá trình nhận dạng và phân tích. Hệ thống cố gắng giữ
	riêng thông tin thành phần, dị nguyên được khai báo và cảnh báo có
	thể chứa vết để kết quả hiển thị rõ ràng hơn.
	
	\Needspace{3\baselineskip}
	\subsection{Nhận dạng và trích xuất văn bản từ hình ảnh}
	
	Nhận dạng ký tự quang học (Optical Character Recognition -- OCR)
	là công nghệ cho phép chuyển đổi chữ viết xuất hiện trong hình ảnh
	thành dữ liệu văn bản có thể xử lý bằng máy tính. Công nghệ này
	được ứng dụng trong việc số hóa tài liệu, nhận dạng thông tin trên
	hóa đơn và trích xuất nội dung từ hình ảnh chứa văn bản.
	
	Đối với bài toán đọc nhãn thực phẩm, quá trình nhận dạng văn bản
	giúp chuyển nội dung được in trên bao bì thành dữ liệu phục vụ
	phân tích. Văn bản sau khi trích xuất có thể chứa tên sản phẩm,
	danh sách thành phần và các thông tin cảnh báo dị nguyên.
	
	Trong đồ án, hệ thống sử dụng Gemini API để hỗ trợ đọc và trích
	xuất thông tin từ hình ảnh nhãn thực phẩm. Thay vì chỉ nhận dạng
	các ký tự riêng lẻ, mô hình AI còn hỗ trợ phân tích nội dung và
	trả về những thông tin có cấu trúc để hệ thống tiếp tục xử lý.
	
	Quá trình xử lý hình ảnh và trích xuất thông tin có thể được mô tả
	qua các bước:
	
	\begin{itemize}
		\item Tiếp nhận hình ảnh nhãn thực phẩm do người dùng tải lên.
		\item Kiểm tra hình ảnh đầu vào và chuẩn bị dữ liệu để xử lý.
		\item Gửi hình ảnh đến dịch vụ Gemini API.
		\item Nhận nội dung văn bản và thông tin sản phẩm được trích xuất.
		\item Kiểm tra dữ liệu trả về và chuyển sang bước phân tích thành phần.
	\end{itemize}
	
	Chất lượng hình ảnh có ảnh hưởng trực tiếp đến kết quả nhận dạng.
	Ảnh bị mờ, thiếu sáng, nghiêng hoặc có kích thước chữ quá nhỏ có
	thể khiến mô hình đọc sai hoặc bỏ sót một số thông tin. Ngoài ra,
	cách bố trí nội dung trên bao bì và ngôn ngữ sử dụng cũng có thể
	ảnh hưởng đến kết quả phân tích.
	
	Do đó, kết quả trích xuất cần được kiểm tra trước khi sử dụng trong
	các bước đối chiếu và tạo cảnh báo.
	
	\Needspace{3\baselineskip}
	\subsection{Tiền xử lý và chuẩn hóa văn bản}
	
	Sau khi trích xuất nội dung từ hình ảnh, dữ liệu văn bản có thể
	chứa khoảng trắng không cần thiết, cách viết hoa và viết thường
	không đồng nhất hoặc những cách diễn đạt khác nhau của cùng một
	thành phần. Vì vậy, tiền xử lý và chuẩn hóa văn bản là bước cần
	thiết để hỗ trợ quá trình phân tích.
	
	Trong hệ thống Allergy Warning, quá trình chuẩn hóa có thể bao gồm:
	
	\begin{itemize}
		\item Loại bỏ khoảng trắng dư thừa và chuẩn hóa định dạng văn bản.
		\item Chuẩn hóa chữ hoa và chữ thường khi thực hiện so sánh.
		\item Tách nội dung danh sách thành phần thành các nguyên liệu riêng biệt khi dữ liệu đầu vào cho phép.
		\item Chuẩn hóa cách biểu diễn tên thành phần để hỗ trợ việc đối chiếu với dữ liệu dị nguyên.
		\item Giữ lại nội dung gốc để phục vụ việc kiểm tra và hiển thị thông tin cho người dùng.
	\end{itemize}
	
	Ví dụ, một thành phần có thể được biểu diễn dưới dạng
	``Whey Powder (Milk)'' hoặc ``whey powder (milk)''. Khi thực hiện
	so sánh, hệ thống có thể chuẩn hóa chữ hoa và chữ thường để hạn chế
	việc bỏ sót do khác biệt về định dạng.
	
	Bên cạnh đó, một số nguyên liệu có thể được gọi bằng nhiều tên
	khác nhau. Vì vậy, quá trình chuẩn hóa cần kết hợp với dữ liệu
	từ khóa hoặc thông tin ánh xạ để hỗ trợ nhận diện thành phần.
	
	Việc chuẩn hóa không đồng nghĩa với việc tự động kết luận một
	nguyên liệu là chất gây dị ứng. Đây là bước chuẩn bị dữ liệu để
	các chức năng phân tích và đối chiếu được thực hiện nhất quán hơn.
	
	\Needspace{3\baselineskip}
	\subsection{Phân tích thành phần và đối chiếu dữ liệu dị nguyên}
	
	Đối chiếu dữ liệu là quá trình so sánh thông tin được trích xuất
	từ nhãn thực phẩm với danh sách dữ liệu dị nguyên đã được lưu trữ.
	Đây là một trong những bước quan trọng của hệ thống Allergy Warning.
	
	Hệ thống sử dụng dữ liệu dị nguyên làm cơ sở tham chiếu để xác
	định những thành phần có khả năng liên quan đến dị ứng. Quá trình
	đối chiếu được thực hiện thông qua các bước cơ bản:
	
	\begin{itemize}
		\item Tiếp nhận danh sách thành phần đã được trích xuất.
		\item Chuẩn hóa nội dung và xác định các tên nguyên liệu.
		\item Đối chiếu tên thành phần với dữ liệu dị nguyên.
		\item Xác định nhóm dị nguyên tương ứng với thành phần được nhận diện.
		\item Tổng hợp kết quả và chuyển sang bước so sánh với hồ sơ dị ứng của người dùng.
	\end{itemize}
	
	Một thành phần có thể được biểu diễn bằng nhiều tên gọi khác nhau.
	Ví dụ, thông tin ``whey powder (milk)'' có thể được nhận diện là
	thành phần liên quan đến sữa. Vì vậy, hệ thống cần xem xét cả tên
	thành phần và các thông tin liên quan thay vì chỉ tìm kiếm một
	từ khóa duy nhất.
	
	Sau khi nhận diện thành phần, hệ thống tiếp tục đối chiếu kết quả
	với hồ sơ dị ứng của người dùng. Nếu dị nguyên được phát hiện
	trùng với dị nguyên mà người dùng đã khai báo, hệ thống hiển thị
	cảnh báo tương ứng.
	
	Cần phân biệt giữa hai kết quả:
	
	\begin{itemize}
		\item Dị nguyên được nhận diện trong sản phẩm: Những nhóm chất có khả năng gây dị ứng được hệ thống xác định từ nội dung thành phần hoặc dữ liệu phân tích.
		\item Dị nguyên trùng với hồ sơ người dùng: Những dị nguyên được nhận diện trong sản phẩm đồng thời nằm trong danh sách dị ứng mà người dùng đã khai báo.
	\end{itemize}
	
	Phương pháp đối chiếu giúp kết quả cảnh báo có căn cứ từ dữ liệu
	được lưu trữ. Tuy nhiên, độ chính xác vẫn phụ thuộc vào chất lượng
	ảnh, nội dung trích xuất, mức độ đầy đủ của dữ liệu dị nguyên và
	cách biểu diễn tên nguyên liệu.
	
	\Needspace{3\baselineskip}
	\subsection{Ứng dụng trí tuệ nhân tạo trong phân tích nhãn thực phẩm}
	
	Trí tuệ nhân tạo (Artificial Intelligence -- AI) là lĩnh vực
	nghiên cứu các phương pháp giúp máy tính thực hiện những nhiệm vụ
	đòi hỏi khả năng phân tích và xử lý thông tin. Trong đồ án
	Allergy Warning, AI được sử dụng để hỗ trợ đọc hình ảnh nhãn
	thực phẩm và phân tích nội dung thành phần.
	
	Hệ thống tích hợp Gemini API để gửi hình ảnh và yêu cầu mô hình
	trích xuất thông tin cần thiết. Kết quả trả về được sử dụng để
	xác định tên sản phẩm, thương hiệu, danh sách thành phần và các
	thông tin cảnh báo dị nguyên khi nội dung trên nhãn cho phép
	nhận diện.
	
	AI có thể hỗ trợ xử lý những trường hợp văn bản có cách diễn đạt
	khác nhau hoặc nội dung được trình bày không đồng nhất. Tuy nhiên,
	kết quả do mô hình tạo ra vẫn có thể chứa thông tin thiếu chính
	xác, bỏ sót thành phần hoặc diễn giải chưa phù hợp với nội dung
	thực tế trên bao bì.
	
	Vì vậy, trong hệ thống cần phân biệt giữa kết quả AI trích xuất
	và kết quả đối chiếu với dữ liệu dị nguyên đã lưu trữ. Hai bước
	này có vai trò bổ trợ cho nhau:
	
	\begin{itemize}
		\item Phân tích bằng AI: Hỗ trợ đọc hình ảnh, trích xuất và diễn giải thông tin trên nhãn thực phẩm.
		\item Đối chiếu dữ liệu: Sử dụng danh sách dị nguyên đã được lưu trữ để kiểm tra các thành phần được nhận diện.
		\item So sánh hồ sơ dị ứng: Kiểm tra những dị nguyên được phát hiện có trùng với thông tin người dùng đã khai báo hay không.
	\end{itemize}
	
	Việc kết hợp AI với dữ liệu dị nguyên giúp hệ thống thực hiện
	quá trình phân tích từ hình ảnh đến kết quả cảnh báo. Tuy nhiên,
	AI không thay thế hoàn toàn việc kiểm tra thông tin trên bao bì
	và không thể bảo đảm phát hiện tất cả các chất gây dị ứng.
	
	Do đó, kết quả của Allergy Warning chỉ có tính chất hỗ trợ
	tham khảo. Người dùng cần kiểm tra lại thông tin trên sản phẩm
	thực tế, đặc biệt trong trường hợp có tiền sử dị ứng nghiêm trọng.
	
	% ---------------------------------------------------------------------
	% 2.2. CÔNG NGHỆ VÀ CÔNG CỤ SỬ DỤNG
	% ---------------------------------------------------------------------
	
	\Needspace{4\baselineskip}
	\section{Công nghệ và công cụ sử dụng}
	
	Để xây dựng hệ thống Allergy Warning, nhóm sử dụng các công nghệ
	phục vụ việc phát triển ứng dụng web, xử lý hình ảnh, phân tích
	thành phần, quản lý dữ liệu và lưu trữ hình ảnh trực tuyến.
	
	Các công nghệ được lựa chọn nhằm đáp ứng những chức năng chính
	của hệ thống, bao gồm quản lý tài khoản, quản lý hồ sơ dị ứng,
	quét nhãn thực phẩm, phân tích dị nguyên và lưu trữ lịch sử quét.
	
	\Needspace{3\baselineskip}
	\subsection{Ngôn ngữ lập trình Python}
	
	Python là ngôn ngữ lập trình bậc cao, có cú pháp tương đối dễ đọc
	và hỗ trợ nhiều thư viện phục vụ phát triển ứng dụng web, xử lý
	dữ liệu và tích hợp các dịch vụ bên ngoài.
	
	Trong đồ án, Python được sử dụng để xây dựng phần xử lý phía máy
	chủ của hệ thống. Ngôn ngữ này đảm nhiệm việc tiếp nhận yêu cầu
	từ giao diện, xử lý dữ liệu đầu vào, gọi dịch vụ AI, thực hiện
	đối chiếu dị nguyên và kết nối cơ sở dữ liệu.
	
	Việc sử dụng Python giúp nhóm tổ chức các chức năng nghiệp vụ
	thành những thành phần riêng biệt, thuận tiện cho quá trình
	phát triển, kiểm thử và bảo trì hệ thống.
	
	\Needspace{3\baselineskip}
	\subsection{Framework Flask}
	
	Flask là một framework web được xây dựng trên nền tảng Python,
	hỗ trợ phát triển ứng dụng web và xây dựng các tuyến xử lý
	yêu cầu từ phía người dùng.
	
	Trong Allergy Warning, Flask được sử dụng để xây dựng máy chủ
	ứng dụng và kết nối các chức năng của hệ thống. Khi người dùng
	thực hiện một thao tác trên giao diện, yêu cầu được gửi đến
	máy chủ Flask để xử lý và trả về kết quả phù hợp.
	
	Các nhiệm vụ chính của Flask trong hệ thống gồm:
	
	\begin{itemize}
		\item Xử lý các yêu cầu đăng ký, đăng nhập và quản lý tài khoản.
		\item Tiếp nhận hình ảnh nhãn thực phẩm do người dùng tải lên.
		\item Gọi dịch vụ Gemini API để hỗ trợ phân tích hình ảnh.
		\item Xử lý dữ liệu thành phần và đối chiếu dị nguyên.
		\item Kết nối cơ sở dữ liệu để lưu trữ và truy xuất thông tin.
		\item Trả kết quả phân tích và lịch sử quét về giao diện web.
	\end{itemize}
	
	Flask đóng vai trò trung tâm trong việc điều phối các thành phần
	xử lý của hệ thống, giúp các chức năng được kết nối thành một
	quy trình thống nhất.
	
	\Needspace{3\baselineskip}
	\subsection{Gemini API}
	
	Gemini API là giao diện lập trình ứng dụng cho phép hệ thống
	tích hợp mô hình trí tuệ nhân tạo Gemini vào ứng dụng của mình.
	Dịch vụ này hỗ trợ tiếp nhận dữ liệu và thực hiện các yêu cầu
	phân tích theo nội dung được gửi đến.
	
	Trong đồ án, Gemini API được sử dụng để hỗ trợ đọc và phân tích
	hình ảnh nhãn thực phẩm. Máy chủ gửi hình ảnh cùng yêu cầu xử lý
	đến dịch vụ AI. Sau đó, hệ thống tiếp nhận kết quả và xử lý
	thông tin trả về.
	
	Các nhiệm vụ chính của Gemini API trong hệ thống gồm:
	
	\begin{itemize}
		\item Hỗ trợ nhận diện nội dung văn bản trên ảnh nhãn thực phẩm.
		\item Trích xuất tên sản phẩm và thương hiệu khi có thông tin.
		\item Trích xuất danh sách thành phần nguyên liệu.
		\item Hỗ trợ nhận diện thông tin dị nguyên được khai báo riêng.
		\item Hỗ trợ nhận diện nội dung cảnh báo có thể chứa vết.
	\end{itemize}
	
	Dữ liệu do Gemini API trả về được máy chủ tiếp nhận và kiểm tra
	trước khi sử dụng cho các bước phân tích tiếp theo. Việc sử dụng
	dịch vụ AI giúp hệ thống xử lý hình ảnh nhãn thực phẩm mà không
	phải xây dựng riêng một mô hình nhận dạng hình ảnh từ đầu.
	
	Tuy nhiên, kết quả phụ thuộc vào chất lượng hình ảnh, nội dung
	yêu cầu gửi đến mô hình và khả năng xử lý của dịch vụ. Vì vậy,
	hệ thống cần có bước kiểm tra dữ liệu trả về và không xem kết
	quả AI là thông tin tuyệt đối chính xác.
	
	\Needspace{3\baselineskip}
	\subsection{TiDB Cloud và cơ sở dữ liệu MySQL}
	
	Cơ sở dữ liệu là thành phần dùng để lưu trữ, quản